\chapter{Giới thiệu}
\section{Đặt vấn đề}
Thực trạng hiện nay cho thấy, mặc dù nhu cầu về chăm sóc sức khỏe và làm đẹp (wellness) đang tăng trưởng mạnh mẽ, cả người dùng và các nhà cung cấp dịch vụ vẫn đang đối mặt với nhiều rào cản đáng kể trong quá trình kết nối và tương tác.
\begin{itemize}
  \item \textbf{Đối với người dùng.}
    \begin{enumerate}
        \item \textbf{Quá tải thông tin và thiếu tin cậy:} Người dùng thường bị "ngợp" trong một ma trận thông tin từ nhiều nguồn khác nhau như mạng xã hội, website riêng lẻ, quảng cáo... Họ gặp khó khăn trong việc xác thực chất lượng, tay nghề của chuyên gia và độ tin cậy của một cơ sở trước khi đưa ra quyết định. Các đánh giá (reviews) thường tản mạn, dễ bị thao túng, không tạo được sự tin tưởng tuyệt đối.
        \item \textbf{Khó khăn trong việc lựa chọn dịch vụ phù hợp:}  Khi gặp một vấn đề sức khỏe cụ thể (ví dụ: đau mỏi vai gáy, căng thẳng), người dùng không biết nên chọn dịch vụ nào là tối ưu nhất: spa thư giãn, y học cổ truyền (châm cứu, bấm huyệt), hay vật lý trị liệu? Họ thiếu một công cụ tư vấn ban đầu, khách quan để đưa ra lựa chọn đúng đắn, thay vì chỉ dựa vào cảm tính hoặc quảng cáo.
        \item \textbf{Quy trình đặt lịch thủ công và thiếu linh hoạt:} Việc đặt lịch hẹn vẫn còn phụ thuộc nhiều vào các phương thức truyền thống như gọi điện, nhắn tin, dẫn đến nhiều bất tiện. Người dùng không có cái nhìn tổng quan về các khung giờ trống, khó thay đổi lịch hẹn và quy trình thanh toán, xác nhận thường không đồng nhất, thiếu chuyên nghiệp.
    \end{enumerate}
  \item \textbf{Đối với nhà cung cấp dịch vụ (Các cơ sở spa, phòng khám...):
}
    \begin{enumerate}
        \item \textbf{Khó khăn trong việc tiếp cận đúng khách hàng tiềm năng:} Các doanh nghiệp vừa và nhỏ, dù có chất lượng dịch vụ tốt, nhưng lại thiếu ngân sách và kỹ năng marketing kỹ thuật số để cạnh tranh, thu hút khách hàng mới một cách hiệu quả trong một thị trường ngày càng sôi động.
        \item \textbf{Quản lý vận hành rời rạc, tốn nhiều nguồn lực:} Việc quản lý lịch hẹn của khách hàng qua nhiều kênh (điện thoại, Zalo, Facebook, sổ sách) rất dễ gây ra sai sót, trùng lặp, bỏ lỡ khách hàng và tốn thời gian của nhân viên. Họ thiếu một hệ thống tập trung để tối ưu hóa quy trình này.
        \item \textbf{Thách thức trong việc xây dựng thương hiệu và uy tín số:} Việc xây dựng lòng tin với khách hàng trên không gian mạng là một quá trình tốn kém và lâu dài. Các nhà cung cấp cần một nền tảng trung gian uy tín để bảo chứng cho chất lượng dịch vụ của mình thông qua các đánh giá khách quan và quy trình xác thực minh bạch.
    \end{enumerate}
\end{itemize}
\textbf{Kết luận}: Thực tế này cho thấy thị trường đang thiếu một \textbf{nền tảng số hóa tập trung, thông minh và đáng tin cậy}, có khả năng giải quyết đồng thời các vấn đề của cả hai phía: đơn giản hóa quá trình tìm kiếm, lựa chọn cho người dùng và tối ưu hóa vận hành, tiếp thị cho nhà cung cấp.

\section{Mục tiêu đồ án}
Đề tài được thực hiện nhằm đạt được các mục tiêu chính và mục tiêu cụ thể sau:

\begin{enumerate}
    \item \textbf{Mục tiêu chính:}
        \begin{itemize}
            \item \textbf{Xây dựng một nền tảng di động toàn diện:} Kết nối người dùng với các nhà cung cấp dịch vụ chăm sóc sức khỏe và sắc đẹp, tạo ra một hệ sinh thái minh bạch và hiệu quả.
            \item \textbf{Ứng dụng Trí tuệ Nhân tạo (AI):} Làm cốt lõi để cá nhân hóa trải nghiệm người dùng, đưa ra các gợi ý dịch vụ thông minh, từ đó nâng cao sự hài lòng và giá trị mang lại.
            \item \textbf{Tối ưu hóa quy trình vận hành:} Cho các đối tác cung cấp dịch vụ thông qua các công cụ quản lý kỹ thuật số, giúp họ tăng trưởng doanh thu và xây dựng thương hiệu bền vững.
        \end{itemize}

    \item \textbf{Mục tiêu cụ thể:}
        \begin{itemize}
            \item \textbf{Phân tích, thiết kế và triển khai ứng dụng di động cho hai đối tượng:} khách hàng và nhà cung cấp dịch vụ (dưới dạng app riêng hoặc module riêng trong cùng một app).
            \item \textbf{Phát triển Chatbot AI:} Tích hợp chatbot có khả năng xử lý ngôn ngữ tự nhiên (NLP) để phân tích nhu cầu, triệu chứng ban đầu của người dùng và đưa ra gợi ý về các loại hình dịch vụ phù hợp nhất với nhu cầu và vị trí người dùng.
            \item \textbf{Xây dựng hệ thống Đặt lịch - Thanh toán - Đánh giá:} Triển khai module cho phép người dùng xem lịch trống, đặt hẹn, thanh toán trực tuyến và để lại đánh giá, nhận xét sau khi sử dụng dịch vụ.
            \item \textbf{Thiết lập quy trình xác thực đối tác:} Xây dựng bộ tiêu chí và quy trình online để kiểm duyệt, xác thực thông tin, giấy phép và chất lượng của các nhà cung cấp trước khi họ tham gia hệ thống.
            \item \textbf{Thiết kế cơ sở dữ liệu tập trung:} Xây dựng kiến trúc CSDL để quản lý hiệu quả và an toàn thông tin người dùng, đối tác, dịch vụ, lịch hẹn và các giao dịch liên quan.
        \end{itemize}
\end{enumerate}

